\section{A complete metric with nonnegative Ricci curvature}
\label{sec:geometry}

Fix spectral data in the sense specified at the end of
Section~\ref{sec:counting}.
The geometric properties we need are \(0<a_*<1<q_*\), the strict
inequality \eqref{eq:link-margin} at \(a=a_*\) for every
\(q\in[1,q_*]\), and the fixed cutoff condition
\eqref{eq:cutoff-smallness}.
Fix also the reference complex structure \(J_0\) used for the dwell
metric.
We describe a metric construction that works for any prescribed finite
family of pulse controls.  Its curvature estimates will be uniform
over all subsequent choices of those controls.
Slowly varying links with fixed volume form also appear in
Colding--Naber~\cite[Lemma~2.1]{ColdingNaber}. We give the curvature
calculation for the present family in full.

The radial cutoff and link parameters were fixed before \(k,L\) and
before any pulse family. Condition~\eqref{eq:cutoff-smallness} will give
\(b=a'/a\geq-1/4\) below. Even very large bounds for a fixed finite
pulse family are handled solely by increasing the starting index \(j_0\).
No spectral parameter is changed in response to the controls.

\paragraph{Pulse data and the period schedule.}
Let \(K\) be a compact ball in a finite-dimensional real parameter
space.
Choose finitely many disjoint open subintervals
\(I_1,\ldots,I_R\) of \((0,1)\), with disjoint closures, smooth
functions \(\beta_\nu\) compactly supported in \(I_\nu\), and
orthogonal complex structures \(J_\nu\).
For each \(\nu\), let \(c_\nu(h)\) be smooth on a neighborhood of \(K\).
Set
\[
 z(\tau,h)=\sum_{\nu=1}^R c_\nu(h)\beta_\nu(\tau).
\]
The amplitudes \(c_\nu\) may have either sign.
The functions, complex structures, and compact set \(K\) are fixed
once and for all.
In applications the bumps can be normalized to have integral \(1\),
so their amplitudes specify the flow parameters in
Corollary~\ref{prop:word}.

For \(j\geq1\), define
\begin{equation}\label{eq:periods}
 T_j=j^6,\qquad L_j=j^7,\qquad t_1=10,\qquad
 t_{j+1}=t_j+L_j+4T_j.
\end{equation}
Choose a starting index \(j_0\), to be made large below.
Put \(q(t)=1\) for \(t\leq t_{j_0}\).
For each \(j\geq j_0\), choose a parameter \(h_j\in K\) and divide
\([t_j,t_{j+1}]\) into the following five successive intervals:
\begin{enumerate}
 \item a pulse of duration \(T_j\), on which
 \[
  q(t)=1+T_j^{-1}z\bigl((t-t_j)/T_j,h_j\bigr);
 \]
 on the support of the \(\nu\)-th bump the complex structure is
 \(J_\nu\);
 \item an upward ramp of duration \(T_j\), taking \(q\) from \(1\)
 to \(q_*\);
 \item a dwell interval of duration \(L_j\), on which \(q=q_*\);
 \item a downward ramp of duration \(T_j\), taking \(q\) from
 \(q_*\) to \(1\);
 \item a round rest of duration \(T_j\), on which \(q=1\).
\end{enumerate}
Use \(J_0\) throughout the last four intervals.
Both ramps use a fixed smooth nondecreasing profile
\(\rho:[0,1]\to[0,1]\), equal to \(0\) near \(0\) and \(1\) near
\(1\): their respective formulas in their own rescaled time are
\(1+(q_*-1)\rho\) and \(q_*-(q_*-1)\rho\).
The last four intervals form the diagonal leg of a period, as
illustrated in Figure~\ref{fig:period}.

\begin{figure}[!ht]
\centering
\setlength{\unitlength}{1mm}
\begin{picture}(152,41)
 \put(0,17){\framebox(25,16){\shortstack{\small Pulse\\$q=1+z/T_j$}}}
 \put(25,17){\framebox(24,16){\shortstack{\small Increase\\$1\to q_*$}}}
 \put(49,17){\framebox(55,16){\shortstack{\small Dwell\\$q=q_*$}}}
 \put(104,17){\framebox(24,16){\shortstack{\small Decrease\\$q_*\to1$}}}
 \put(128,17){\framebox(24,16){\shortstack{\small Round rest\\$q=1$}}}
 \put(12.5,12){\makebox(0,0){$T_j$}}
 \put(37,12){\makebox(0,0){$T_j$}}
 \put(76.5,12){\makebox(0,0){$L_j$}}
 \put(116,12){\makebox(0,0){$T_j$}}
 \put(140,12){\makebox(0,0){$T_j$}}
 \put(25,6){\line(1,0){127}}
 \put(25,6){\line(0,1){2}}
 \put(152,6){\line(0,1){2}}
 \put(88.5,2){\makebox(0,0){\small Fixed diagonal leg: $J=J_0$}}
 \put(0,38){\makebox(0,0)[l]{$t_j$}}
 \put(152,38){\makebox(0,0)[r]{$t_{j+1}$}}
\end{picture}
\caption{A period in logarithmic radius, with $T_j=j^6$ and $L_j=j^7$.
The pulse contains finitely many weak bumps with round gaps between them.
The four subsequent pieces form the fixed diagonal leg. Horizontal lengths
are schematic.}
\label{fig:period}
\end{figure}


The notation \(J(t)\) is only a label for the complex structure used
in the metric at time \(t\).  It need not be continuous.
Labels are changed only in open intervals where \(q=1\), and
\(G(a,1,J)=a^2g_0\) is independent of \(J\) there.
The supported bumps provide such round gaps before, between, and
after their supports.

\begin{proposition}\label{prop:metric}
For the fixed spectral data above, there are a smooth function
\(a:\R\to(0,1]\) and an index \(j_{\rm geom}\) such that, for every
\(j_0\geq j_{\rm geom}\), every sequence
\((h_j)_{j\geq j_0}\) in \(K\), used in the preceding scheme, defines
a smooth complete metric on \(\R^{m+1}\) with nonnegative Ricci
curvature:
\[
 g=dr^2+r^2\gamma(\log r),\qquad
 \gamma(t)=G(a(t),q(t),J(t)).
\]
The function \(a\) is identically \(1\) for \(t\leq2\), decreases to
\(a_*\), and, for some \(\mu>0\), satisfies
\begin{equation}\label{eq:scale-tail}
 b(t):=\frac{a'(t)}{a(t)}
       =-\mu(1+t)^{-5/4}\quad(t\geq4),
 \qquad
 \log\frac{a(t)}{a_*}=4\mu(1+t)^{-1/4}\quad(t\geq4).
\end{equation}
The metric is Euclidean near the origin, and its distance from the
origin is exactly \(r\).
Its volume form is
\[
 d\operatorname{vol}_g
   =r^m a(\log r)^m\,dr\,d\operatorname{vol}_{g_0}.
\]
The function \(a\) and the starting index \(j_0\) can be chosen
independently of the sequence of control parameters.
\end{proposition}

\begin{proof}
\textbf{The radial scale.}
Use the fixed smooth nondecreasing cutoff \(\chi\) from
Section~\ref{sec:counting}.
For \(t>2\), put
\[
 b(t)=-\mu\chi(t)(1+t)^{-5/4},\qquad
 a(t)=\exp\left(\int_2^t b(u)\,du\right),
\]
and set \(b=0\), \(a=1\) for \(t\leq2\).
The integral of \(\chi(t)(1+t)^{-5/4}\) over \((2,\infty)\) is
positive and finite.  Thus there is a unique \(\mu>0\) for which
\(a(t)\to a_*\); moreover \(\mu\to0\) as \(a_*\to1\).
Condition~\eqref{eq:cutoff-smallness} gives \(b\geq-1/4\).
Since \(\chi'\geq0\), for \(t>2\) we have
\[
 b'(t)\leq-\frac{5b(t)}{4(1+t)},\qquad
 b+b'+b^2
 \leq b\left(1-\frac{5}{4(1+t)}+b\right)\leq0.
\]
The last factor is positive for \(t>2\) and \(b\geq-1/4\).
The tail identities \eqref{eq:scale-tail} follow by integration.

\medskip\noindent
\textbf{The slice shape and the mixed Ricci tensor.}
Write \(g_r=r^2\gamma(\log r)\) for the metric on a slice.
A dot will denote differentiation in \(t=\log r\).
Identify the tangent bundles of the slices using the product
coordinates.  Their shape endomorphism is
\[
 S=\frac12 g_r^{-1}\partial_r g_r
   =r^{-1}(\Id+B),\qquad
 B=\frac12\gamma^{-1}\dot\gamma=b\Id+H.
\]
Where a nonround part of a bump or a ramp occurs, \(J\) is fixed.
Direct differentiation of \(G\) gives
\begin{equation}\label{eq:shape-anisotropy}
 H=\frac{\dot q}{2q}\left(P_J-\frac1m\Id\right),\qquad
 \tr H=0,\qquad
 \tr(H^2)=\frac{m-1}{4m}\left(\frac{\dot q}{q}\right)^2.
\end{equation}
On round gaps \(H=0\).  Thus the same identities apply everywhere,
without requiring derivatives of the labels \(J(t)\).

The projection \(P_J\) is also the orthogonal projection for
\(\gamma\) onto the Hopf direction.
Indeed, \(\xi_J\) has constant squared length
\(a^2q^{1-1/m}\) on a fixed slice, and its metric dual is the same
constant times \(\eta_J\).
The flow of \(\xi_J\) preserves both \(g_0\) and \(\eta_J\), so
\(\xi_J\) is a Killing field for \(\gamma\).
A Killing field of constant length has zero divergence and
\(\nabla^\gamma_{\xi_J}\xi_J=0\): the latter identity follows by
pairing the Killing equation with \(\xi_J\).
Equivalently, its unit normalization \(E\) satisfies
\[
 \operatorname{div}_\gamma(E\otimes E^\flat)
    =(\operatorname{div}_\gamma E)E^\flat
      +(\nabla^\gamma_EE)^\flat=0.
\]
Therefore \(\operatorname{div}_\gamma P_J=0\), and
\(\operatorname{div}_\gamma H=0\), because the coefficients in
\eqref{eq:shape-anisotropy} are constant on a slice.
The trace \(\tr S=r^{-1}m(1+b)\) is also spatially constant.
The Codazzi formula consequently makes the mixed Ricci tensor
identically zero.

We record the remaining curvature identities explicitly.
For a metric \(dr^2+g_r\), the connection components involving the
radial variable are
\(\Gamma^r_{ij}=-(g_r)_{ik}S^k{}_j\) and
\(\Gamma^i_{rj}=S^i{}_j\).
They give the normal curvature endomorphism
\(-\partial_rS-S^2\).  Its trace is \(\Ric_g(\partial_r,\partial_r)\).
Tracing the tangential Gauss equation contributes
\(S^2-(\tr S)S\) to the intrinsic slice Ricci endomorphism.
Adding the normal contribution gives
\[
 \Ric_g^\#\big|_{TS^m}
   =\Ric_{g_r}^\#-\partial_rS-(\tr S)S.
\]
Here an index of the tangential tensor is raised using \(g_r\).
The mixed formula is
\(\operatorname{div}_{g_r}S-d(\tr S)\), already shown to vanish.
Substitution of \(S=r^{-1}(\Id+B)\), and
\(\tr B=mb\), yields
\begin{equation}\label{eq:radial-slice-ricci}
 \begin{aligned}
  r^2\Ric_g(\partial_r,\partial_r)
     &=-m(b+\dot b+b^2)-\tr(H^2),\\
  r^2\Ric_g^\#\big|_{TS^m}
     &=\Ric_\gamma^\#-\dot B
       -\bigl(m(1+b)-1\bigr)(\Id+B).
 \end{aligned}
\end{equation}
In particular, the \(H^2\) terms cancel from the tangential identity.

\medskip\noindent
\textbf{Uniform positivity on the periods.}
All constants in the estimates that follow may depend on the fixed
pulse data and \(K\), but not on \(j\) or on any choices of \(h_j\).
On a pulse,
\[
 |q-1|=O(T_j^{-1}),\qquad
 |\dot q|=O(T_j^{-2}),\qquad
 |\ddot q|=O(T_j^{-3}).
\]
On either ramp, \(|\dot q|=O(T_j^{-1})\) and
\(|\ddot q|=O(T_j^{-2})\).
There is no variation of \(q\) during the dwell or rest.
Taking \(j_0\) large makes \(q>0\) throughout all pulses, including
those with negative amplitudes.
The metrics on all slices then lie in a fixed uniformly equivalent
family.  Formula \eqref{eq:shape-anisotropy} gives, throughout period
\(j\),
\[
 \|H\|=O(T_j^{-1}),\qquad
 \|\dot H\|=O(T_j^{-2}),\qquad
 \tr(H^2)=O(T_j^{-2})=O(j^{-12}).
\]
The stronger pulse bound is \(\tr(H^2)=O(T_j^{-4})\).
To obtain the derivative estimate, note that \(P_J\) is independent
of \(t\) where \(q\) varies; on every label-changing gap \(H\)
vanishes identically.

From \eqref{eq:periods},
\[
 t_j\sim\frac{j^8}{8},\qquad
 \sup_{t\in[t_j,t_{j+1}]}\left|\frac{t}{t_j}-1\right|\longrightarrow0.
\]
On the scale tail, an exact calculation gives
\[
 -m(b+\dot b+b^2)
 =m\mu(1+t)^{-5/4}
       \left(1-\frac{5}{4(1+t)}-\mu(1+t)^{-5/4}\right).
\]
It has a positive lower bound of order \(j^{-10}\), uniformly on
period \(j\), for all sufficiently large \(j\).
This dominates the \(O(j^{-12})\) term \(\tr(H^2)\).
The radial component in \eqref{eq:radial-slice-ricci} is therefore
positive once \(j_0\) is large enough.

For the tangential component, the strict inequality
\eqref{eq:link-margin} on the compact interval \([1,q_*]\) implies
that, for some fixed \(\delta>0\),
\[
 \Ric_{G(a_*,q,J_0)}^\#\geq(m-1+2\delta)\Id
       \qquad(1\leq q\leq q_*).
\]
All \(J\) give the same intrinsic Ricci eigenvalues.
Since \(a(t)\to a_*\), this estimate persists with \(2\delta\)
replaced by \(\delta\) along sufficiently late ramps, dwells, and
round rests.
The pulse values satisfy \(q\to1\) uniformly; the strict inequality
at \((a_*,1)\) likewise supplies the same conclusion on late pulses,
after decreasing \(\delta\) if necessary.
This argument allows \(q<1\) on a pulse.
Thus all sufficiently late slices satisfy
\[
 \Ric_\gamma^\#\geq(m-1+\delta)\Id.
\]
On the other hand, \(B=b\Id+H\), \(b,\dot b\to0\), and the displayed
bounds for \(H,\dot H\) show uniformly that
\[
 \dot B+\bigl(m(1+b)-1\bigr)(\Id+B)
       =(m-1)\Id+o(1).
\]
The tangential component of
\eqref{eq:radial-slice-ricci} is consequently positive as well.
All the thresholds used here depend only on the fixed data and
\(K\).  Since the mixed component is zero, these estimates prove
\(\Ric_g\geq0\) on every period for every admitted parameter sequence.

\medskip\noindent
\textbf{The initial region and the global manifold.}
Before \(t_{j_0}\) the metric is round in the angular variable:
\[
 g=dr^2+f(r)^2g_0,\qquad f(r)=r\,a(\log r).
\]
The scale construction gives
\[
 f'(r)=a(1+b)\in(0,1],\qquad
 f''(r)=\frac{a}{r}(b+\dot b+b^2)\leq0.
\]
Its radial and tangential sectional curvatures are respectively
\(-f''/f\) and \((1-(f')^2)/f^2\), as follows from the same normal
curvature and Gauss formulas.  They are nonnegative, proving the
required Ricci condition in this region.

The pulse bumps and ramp profiles are constant or flat at all
junctions, and label changes occur only in round open intervals.
Thus \(\gamma\), and hence \(g\), is smooth for \(r>0\).
Only finitely many periods meet any bounded interval of \(t\), so
the infinite sequence of periods introduces no finite accumulation
of junctions.
For \(t\leq2\), equivalently \(r\leq e^2\), both \(a\) and \(q\)
equal \(1\).  The metric is exactly Euclidean there and extends
smoothly across the origin.
The volume formula follows from Lemma~\ref{lem:berger-metrics}.

A radial path from the origin to a point of radius \(r\) has length
\(r\).  Conversely, the length of any such path is at least the
total variation of its radial coordinate, and hence at least \(r\).
Thus \(d_g(0,x)=r(x)\).
The closed centered balls are the compact sets \(\{r\leq R\}\)
in the underlying \(\R^{m+1}\), so the metric is complete.
The underlying manifold is connected, smooth, and without boundary.
\end{proof}
